

\begin{center}
    \textcolor{red}{\texttt{This section can be removed if the content has already been discussed in a common section.}}
\end{center}

The amplitude for the process $g(p_1)g(p_2) \to H(p_3)H(p_4)$ can be expressed as ${\cal M}=\varepsilon_{1,\mu}\varepsilon_{2,\nu}{\cal M}^{\mu\nu}$, where $\varepsilon_{i,\mu}$ denotes the gluon polarization vector. Thanks to the current conservation relations $p_{1,\mu}{\cal M}^{\mu\nu}=p_{2,\nu}{\cal M}^{\mu\nu}=0$, the amplitude can be decomposed as
\begin{align}
{\cal M}^{\mu\nu}=F_1T_1 ^{\mu\nu}+F_2T_2 ^{\mu\nu}+\Delta_0^{\mu\nu}+\Delta_5^{\mu\nu}
\,,
\end{align}
where $\Delta_5^{\mu\nu}$ contains a single Levi-Civita tensor and appears from two loops onward, and $\Delta_0^{\mu\nu}$ depends on either $p_1^{\mu}$ or $p_2^{\nu}$. These two tensors will vanish when multiplying together the NLO electroweak (EW) and the LO amplitudes. The physical tensors can be expressed as \cite{Glover:1987nx,Plehn:1996wb}
\begin{align}
T_1 ^{\mu\nu} &= g^{\mu\nu}-\frac{p_1^{\nu}\,p_2^{\mu}}{p_1\cdot  p_2}
\,,\label{eq:Ttensors} \\
%&T_2 ^{\mu\nu}= g^{\mu\nu}+\frac{1}{p_{T}^{2}\,(p_1\cdot p_2)}\big[ 2\,(p_1\cdot p_2) \,p_3^{\nu}\,p_3^{\mu}  \nonumber  \\&~~
%- 2\,(p_1\cdot p_3) \,p_3^{\nu}\,p_2^{\mu} -2\,(p_2\cdot p_3) p_3^{\mu}\,p_1^{\nu}+ m_H^2 \,p_1^{\nu}\,p_2^{\mu} \big]  \\
T_2 ^{\mu\nu} &= g^{\mu\nu}+\frac{1}{p_{T}^{2}\,(p_1\cdot p_2)}\left[ 2(p_1\cdot p_2)p_3^{\mu}p_3^{\nu} - 2(p_1\cdot p_3)p_2^{\mu}p_3^{\nu}
%\right.\nonumber\\&~~~~~~~~~~~~~~~~~~~~~~~~~~~~~~~~~~~~~~\left.
-2(p_2\cdot p_3)p_3^{\mu}p_1^{\nu}+ m_h^2 p_2^{\mu}p_1^{\nu} \right]
, \label{eq:Ttensors2}
\end{align}
where $p_{T}=\sqrt{(\hat{t}\hat{u} - m_h^4)/\hat{s}}$ represents the transverse momentum of one of the Higgs with 
\begin{eqnarray}
&&\hat{s}=(p_1+p_2)^2,~\hat{t}=(p_1-p_3)^2,~\hat{u}=(p_2-p_3)^2.\label{st}
\end{eqnarray} 
The gauge-invariant form factors $F_{1,2}$ can be extracted by defining the projectors $P_{1,2}^{\mu\nu}$ in terms of $T_{1,2}^{\mu\nu}$ as
\begin{eqnarray}
P_1^{\mu\nu} &=&+\frac{1}{4}\,\frac{D-2}{D-3} \,T_1^{\mu\nu}
-\frac{1}{4}\,\frac{D-4}{D-3} \,T_2^{\mu\nu}
\,,\label{eq:proj1}\\
P_2^{\mu\nu}&=& -\frac{1}{4}\,\frac{D-4}{D-3} \,T_1^{\mu\nu}
+\frac{1}{4}\,\frac{D-2}{D-3} \,T_2^{\mu\nu}
\,,\label{eq:proj2}
\end{eqnarray}   
such that 
\begin{eqnarray}
F_1=P_1^{\mu\nu} {\cal M}_{\mu\nu}\,,\qquad F_2=P_2^{\mu\nu} {\cal M}_{\mu\nu}\,.
\end{eqnarray}

Real-emission processes, such as $gg \to hh\gamma$, are forbidden by Furry's theorem, restricting NLO EW corrections to purely virtual contributions at the two-loop level.




\begin{center}
    \textcolor{red}{\texttt{A shorter version of this section, with only EW relevant results.}}
\end{center}